\documentclass[11pt,a4paper]{article}
\usepackage{kagan}
\renewcommand{\headlabel}{Roadmap}
\hypersetup{pdftitle={Kagan: Roadmap and Architecture}, pdfauthor={Kagan}}

\begin{document}
\thispagestyle{plain}

\noindent\begin{tikzpicture}
  \path[use as bounding box] (0,0) rectangle (\linewidth,4.2cm);
  \fill[tkNavy, rounded corners=5pt] (0,0) rectangle (\linewidth,4.2cm);
  \fill[tkRed] (0.05\linewidth,0.9cm) rectangle (0.058\linewidth,3.3cm);
  \node[font=\sffamily\bfseries\fontsize{30}{32}\selectfont, text=white, anchor=north west] at (0.09\linewidth,3.45cm) {Kagan};
  \node[font=\sffamily\bfseries\Large, text=white, anchor=north west] at (0.09\linewidth,2.15cm) {Roadmap and architecture};
  \node[font=\sffamily\small, text=white!78!tkNavy, anchor=south west, text width=0.86\linewidth, align=left] at (0.09\linewidth,0.45cm)
    {A UNIX-like operating system for i386, RV32 and AArch64: design decisions, development phases, release milestones and risks};
\end{tikzpicture}

\vspace{0.6em}
\tableofcontents

\section{Introduction}

Kagan is a small UNIX-like operating system written from scratch in C and assembly. Release 1.0 targets 32-bit
x86 PCs: it boots with GRUB, runs preemptive multitasking with a separate address space per process, loads ELF
programs from its own on-disk file system, and provides a POSIX-flavoured system call interface, a C library and a
shell with pipes and redirection. Release 2.0 moves everything that depends on the processor behind one architecture
interface and runs the same kernel, C library and file-system image on 32-bit RISC-V and 64-bit ARM. Every target
runs in QEMU first; real hardware follows once the emulated target is stable.

This document is the entry point to the design documentation. It records the technical decisions that hold for the
whole project, the nineteen development phases and the order between them, the architecture of releases 1.0 and
2.0, the project risks, and the exit criteria of every phase.

\subsection{Design goals}

\begin{description}
  \item[Small and readable.] A monolithic kernel with a modular source tree and a short path from a system call to
    the code that serves it. Assembly only where C cannot work, well under 5\,\% of the code.
  \item[Correct at the protection boundary.] Every user pointer goes through one file, every system call has a
    documented validation contract, and every claim about isolation is backed by a negative test or a fuzzer.
  \item[Portable by construction.] Generic code never sees the CPU; firmware differences end at one boot-information
    structure filled from Multiboot or a devicetree.
  \item[Compatible where it costs nothing.] Linux i386 system call numbers on every architecture, and the MINIX~V3
    on-disk format, which Linux's \texttt{mkfs.minix} and \texttt{fsck.minix} create and check.
  \item[Tested at every step.] Each phase closes with exit criteria that run headless in QEMU and in continuous
    integration.
\end{description}

\subsection{The document set}

Each phase has its own design document in \texttt{docs/pdf/}, built from the LaTeX sources in \texttt{docs/phases/}.
All of them share one structure: an at-a-glance box with the phase's dependencies and deliverables; \emph{Scope and
motivation}; \emph{Background}, the concepts the implementation relies on; \emph{References}, with chapter, section
and page references weighted \must, \should or \could; \emph{Implementation}, a sequence of numbered tasks, each
with an acceptance check marked \emph{Verify}; \emph{Testing and debugging}, a table of failure modes; \emph{Exit
criteria}; review questions; and \emph{Extensions}. A \emph{Prior art} box compares the design with MINIX~3 and
Linux, and in Stage V with xv6-riscv and Linux. The Stage V documents also cite free specifications; their chips are
explained in a second legend at the end of those documents and of this one.

\subsection{Development process}

Phases are implemented in order, because each one builds on everything before it. Work proceeds in small commits,
each of which builds and boots. A task is complete only when its Verify check passes, and a phase is complete only
when all of its exit criteria hold; the next phase does not start while any of them is still open. Design
decisions, hypotheses and debugging notes are recorded in the engineering journal, \texttt{docs/journal.md}.

\section{References}

Seven texts underpin the design. \OSTEP, \OSC and \MOS provide the theory; \OSDI documents a complete small kernel,
MINIX~3, which serves as the main prior-art comparison; \LKD and \LKP show how Linux implements each subsystem today;
and \LOB is a practical guide to x86 bring-up whose chapters map onto Phases 1--14.

\begin{table}[htbp]
\centering\small\renewcommand{\arraystretch}{1.4}
\begin{tabularx}{\linewidth}{@{}l >{\raggedright\arraybackslash}p{0.27\linewidth} >{\raggedright\arraybackslash}X@{}}
\toprule
\tkth{Code} & \tkth{Source} & \tkth{Role in the design} \\
\midrule
\LOB & The little book about OS development & Practical x86 bring-up reference; its 14 chapters follow the build order of Phases 1--14. It uses GRUB Legacy and Bochs; the phase documents give the GRUB 2 and QEMU equivalents. \\
\OSTEP & Operating Systems: Three Easy Pieces & Primary theory reference for every phase. \\
\OSC & Operating System Concepts, 10th ed. & Comprehensive reference; a second treatment of each topic. \\
\MOS & Modern Operating Systems, 4th ed. & Hardware, I/O, clocks, disks, security and operating-system design (Ch.~12). \\
\OSDI & Operating Systems: Design and Implementation, 3rd ed. & MINIX~3, a complete x86 kernel printed line by line: the main prior-art comparison. MINIX is a microkernel; Kagan is monolithic. \\
\LKD & Linux Kernel Development, 3rd ed. & How Linux implements scheduling, interrupts and bottom halves, synchronisation, memory, the VFS and block I/O: the second prior-art comparison. \\
\LKP & Linux Kernel Programming (2021) & Current Linux practice for logging, the buddy and slab allocators, the CFS scheduler and synchronisation. \\
\bottomrule
\end{tabularx}
\caption{The seven reference texts and their role in the design.}
\end{table}

\subsection{Editions}\label{sec:getbooks}

\LOB and \OSTEP are freely available from their authors; the others are published textbooks.
Table~\ref{tab:editions} lists the exact editions the reference tables cite.

\begin{table}[htbp]
\centering\small\renewcommand{\arraystretch}{1.4}
\begin{tabularx}{\linewidth}{@{}l >{\raggedright\arraybackslash}p{0.45\linewidth} >{\raggedright\arraybackslash}X@{}}
\toprule
\tkth{Code} & \tkth{Edition cited} & \tkth{Availability} \\
\midrule
\LOB & Erik Helin, Adam Renberg. \emph{The little book about OS development}, 2015. &
  Free at \url{https://littleosbook.github.io}, as a web page or a PDF. Page numbers refer to the PDF. \\
\OSTEP & Remzi H.\ Arpaci-Dusseau, Andrea C.\ Arpaci-Dusseau. \emph{Operating Systems: Three Easy Pieces}, version 0.90, Arpaci-Dusseau Books, 2015. &
  Free at \url{https://pages.cs.wisc.edu/~remzi/OSTEP/}, one PDF per chapter, in the current version (1.10). \\
\OSC & Abraham Silberschatz, Peter Baer Galvin, Greg Gagne. \emph{Operating System Concepts}, 10th ed., Wiley, 2018. ISBN 978-1-119-32091-3. &
  Wiley. \\
\MOS & Andrew S.\ Tanenbaum, Herbert Bos. \emph{Modern Operating Systems}, 4th ed., Pearson, 2015. ISBN 978-0-13-359162-0. &
  Pearson. \\
\OSDI & Andrew S.\ Tanenbaum, Albert S.\ Woodhull. \emph{Operating Systems: Design and Implementation}, 3rd ed., Pearson Prentice Hall, 2006. ISBN 978-0-13-142938-3. &
  Pearson Prentice Hall (out of print; widely available second-hand). \\
\LKD & Robert Love. \emph{Linux Kernel Development}, 3rd ed., Addison-Wesley, 2010. ISBN 978-0-672-32946-3. &
  Addison-Wesley. \\
\LKP & Kaiwan N.\ Billimoria. \emph{Linux Kernel Programming}, Packt, 2021. ISBN 978-1-78995-343-5. &
  Packt. \\
\bottomrule
\end{tabularx}
\caption{Editions cited in the design documents.}
\label{tab:editions}
\end{table}

\begin{warnbox}[Page numbers belong to these editions]
The page ranges in the reference tables are the \emph{printed} page numbers of the editions in
Table~\ref{tab:editions}. They do not match other editions, and a reflowable e-book has no fixed pages, so each
reference also gives chapter and section numbers. For \OSTEP this is the normal case: each free chapter PDF starts at
page~1, and the current version is paginated differently from version 0.90. Chapters 1--43, and every section the
documents cite, have the same numbers in both versions. One title changed: chapter 23, \emph{The VAX/VMS Virtual
Memory System}, is now \emph{Complete Virtual Memory Systems}, with VAX/VMS in \S23.1.
\end{warnbox}

\subsection{Specifications and other references}

Exact hardware behaviour comes from specifications rather than textbooks. The \textbf{Intel 64 and IA-32
Architectures Software Developer's Manual, Volume 3A} (System Programming Guide) is the authority for protection,
segmentation, paging, interrupts and the TSS, chapters 2--7. The \textbf{Multiboot Specification 0.6.96} defines the
contract between GRUB and the kernel, and the \textbf{NASM manual} the assembler syntax. The \textbf{OSDev Wiki}
(\url{https://wiki.osdev.org}) documents nearly every device used here, with known pitfalls. MIT's \textbf{xv6} is a
useful small UNIX-like comparison kernel.

Stage V leaves x86, so the Intel manual, Multiboot and NASM no longer apply there. Its documents cite four further
free references, each with its own chip. \XVSIX is the xv6 book, RISC-V edition, read with the \texttt{xv6-riscv}
source, the closest small kernel to the Phase 17 port. \RVSPEC is the RISC-V Instruction Set Manual, Volume I
(unprivileged) and Volume II (privileged architecture), plus the RISC-V Supervisor Binary Interface (SBI)
specification. \ARMDOC is the Arm Architecture Reference Manual for A-profile (DDI~0487) together with Arm's
\emph{Learn the architecture} guides. \SPEC covers the Devicetree and Virtio specifications and QEMU's documentation
for its \texttt{virt} boards. The port legend at the end of this document gives full titles and links.

\section{Key technical decisions}\label{sec:decisions}

These decisions hold for the whole project. The first eight follow the reference texts, so the design can be compared
with MINIX~3 and Linux at every step; the last five cover the Stage V ports and follow the free specifications.

{\small\renewcommand{\arraystretch}{1.4}
\begin{xltabular}{\linewidth}{@{}>{\raggedright\arraybackslash}p{0.17\linewidth} >{\raggedright\arraybackslash}p{0.27\linewidth} X@{}}
\toprule
\tkth{Decision} & \tkth{Choice} & \tkth{Why} \\
\midrule
\endfirsthead
\toprule
\tkth{Decision} & \tkth{Choice} & \tkth{Why} \\
\midrule
\endhead
\bottomrule
\caption{Technical decisions for Kagan.}\\
\endlastfoot
Target CPU & IA-32 (i386), 32-bit protected mode, BIOS boot & \LOB, \OSDI (MINIX 3) and the Intel examples in \OSC\,\S9.6 are all 32-bit x86. Two-level paging is the simplest real MMU to learn. A 64-bit x86 port stays a Phase 15 track; the RISC-V and ARM ports are Stage V. \\
Boot loader & GRUB via Multiboot~1; \texttt{qemu -kernel} for quick runs & A custom boot loader adds no value at this scale. Phase 1 implements a 512-byte boot sector to establish the boot path, then GRUB takes over. \\
Languages & C (gnu11) plus NASM (Intel syntax) & Assembly only where C cannot work: the entry point, descriptor table loads, interrupt stubs, context switch and entering user mode. That is under 5\,\% of the code. \\
Toolchain & \texttt{i686-elf} cross GCC and binutils, GNU Make, Git & A cross compiler never assumes a Linux target, which prevents a family of confusing bugs (Phase 0). \\
Emulator, debugger & QEMU and GDB (remote stub); Bochs is optional & QEMU is fast, scriptable, and has a monitor for inspecting registers and page tables. \LOB uses Bochs, whose debugger is still handy for descriptor-table bugs. \\
Kernel structure & Monolithic, modular source tree & Simplest to get right for a first kernel. MINIX's microkernel is studied as a contrast, with an optional experiment in Phase 15. \\
On-disk file system & MINIX V3 layout with 1\,KiB blocks, named KaganFS & The format described in \OSDI\,\S5.6, with source. Linux's \texttt{mkfs.minix -3} and \texttt{fsck.minix} create and check the images, and Linux can even mount them. \\
Testing & In-kernel \texttt{ktest} suites, headless QEMU, \texttt{isa-debug-exit}; host tests and parser fuzzing under AddressSanitizer and UBSan; continuous integration & \texttt{make test} boots, runs every test and exits with a pass or fail code, so regressions show up at once (Phase 7). CI runs it on every push; negative tests (``evil programs'') and a system call fuzzer check the protection boundary from Phases 10 and 11 on. \\
Port targets & RV32 (\texttt{rv32imac}) with Sv32 paging on QEMU \texttt{virt}, then optionally the custom Timur-RV32IMC core; AArch64 (ARMv8-A) at EL1 with 39-bit virtual addresses on QEMU \texttt{virt}, then optionally a Raspberry Pi 4 &
  Sv32 is two-level 4\,KiB paging like x86, so the RISC-V kernel can stay at \texttt{0xC0000000}. AArch64 then changes the instruction set and the word size together: 4\,KiB pages, three-level tables, the kernel in the \texttt{TTBR1} high half (Phases 17, 18). \\
Port firmware and discovery & OpenSBI in M-mode (the kernel runs in S-mode); QEMU or the Pi firmware on ARM; a flattened devicetree instead of Multiboot &
  The firmware hands over a devicetree that says where RAM and devices are. \texttt{lib/fdt.c} reads it into \texttt{struct bootinfo}, so generic code never knows which firmware booted it (Phase 16). \\
System calls on the ports & The same Linux i386 numbers everywhere; \texttt{ecall} with the number in \texttt{a7}, \texttt{svc \#0} with it in \texttt{x8} &
  One table and one libc for three architectures; only the trap instruction and the registers change. Linux itself uses a different, generic table on RISC-V and arm64 (\texttt{write} is 64 and there is no \texttt{fork}), so its headers are not a source for these numbers. \\
Port toolchain & \texttt{riscv32-elf} and \texttt{aarch64-elf} cross GCC and binutils; GNU \texttt{as} for port assembly & Built exactly like \texttt{i686-elf} in Phase 0, with the same versions (Task 16.1). NASM only knows x86, so it stays for the i386 code; port assembly lives in preprocessed \texttt{.S} files. \\
Port devices & 16550 or PL011 console, PLIC or GICv2, SBI or Arm generic timer, virtio-blk disk, \texttt{sifive\_test} or semihosting for \texttt{make test} &
  The \texttt{virt} boards have no PS/2 keyboard, VGA text screen or IDE disk by default. One virtio-blk driver serves both ports. All three targets are little-endian and KaganFS uses fixed-width fields, so the same disk image works everywhere. \\
\end{xltabular}}

\begin{decisionbox}{Why monolithic rather than a microkernel?}
In a monolithic kernel, drivers, the file system and the scheduler all run in ring 0 and call each other as plain C
functions. In a microkernel like MINIX 3 they run as separate user-space processes that talk by message passing
(\OSDI\,\S1.5.5, \S2.5.1; \MOS\,\S1.7). Microkernels isolate faults better, but they require IPC, servers and a
protocol for every operation before anything works. Plain function calls keep the path from a design to running code
short and make the kernel easier to reason about. Every phase still compares the design with MINIX, and Phase 15
specifies an experiment that moves one driver into a user-space server to measure the trade-off.
\end{decisionbox}

\begin{decisionbox}{Security starts where its subject starts}
The security goal for release 1.0 is concrete: no unprivileged program and
no malformed disk image may crash or take over the kernel, and every claim of that kind is backed by a test. Each
practice begins in the phase where the thing it protects first exists. Phase~0 writes a short threat-model stub and
verifies the toolchain downloads; Phase~2 turns every warning into an error; Phase~6 makes kernel code read-only;
Phase~7 adds continuous integration and a build that traps on undefined behaviour; Phase~10 writes the real threat
model, routes every user pointer through one small file and starts a table of how each system call checks its
arguments; Phase~11 starts the system call fuzzer; Phase~13 treats every on-disk structure as untrusted input.
Phase~15 then consolidates: it finishes the audit, adds users and permissions, and ships a release whose build can
be reproduced and checked. The model, its trust boundaries and what is out of scope (physical attackers, the boot
chain, networks, malicious hardware, side channels) live in \texttt{docs/security/threat-model.md}.
\end{decisionbox}

\begin{decisionbox}{When the references disagree}
The reference texts span fifteen years, from \OSDI (2006) to \LKP (2021), and sometimes describe different ways to
do the same job. The rule: adopt the newer practice when it also suits Kagan, and keep the older one when it is
simpler and the newer one only pays off in a much larger kernel. So Kagan
\textbf{adopts} printk-style log levels and a \code{dmesg} buffer (Phase~2), top and bottom halves for interrupts
(Phase~4), the 0/\code{-E} and \code{ERR_PTR} error convention (Phase~7), saturating \code{refcount_t} counters and
a run-time lock-order check (Phase~9). It \textbf{stages} the buddy page allocator after the bitmap (Phase~5) and
slab caches under \code{kmalloc} after the first-fit heap (Phase~7), so the simple version is in place and tested first. It
offers a CFS-style \textbf{fair scheduler} instead of MLFQ as the optional scheduler (Phase~8). And it
\textbf{keeps} the MINIX-style buffer cache rather than Linux's page cache (Phase~12), and leaves loadable modules
to the optional Track~G, because a small monolithic kernel gains little from either.
\end{decisionbox}

\begin{decisionbox}{Why port after 1.0, and why RISC-V before ARM?}
A complete system on one architecture comes first, and x86 is the architecture of the main references. Porting
afterwards is not wasted effort: every place where x86 leaked into generic code shows up when it is compiled for
another machine, and fixing those leaks is what makes the generic code genuinely generic. RISC-V
comes first because RV32 keeps the 32-bit word size, Sv32 paging has the same two-level, 4\,KiB shape as x86, so
the kernel can stay at \texttt{0xC0000000}, and the privileged architecture is small enough to specify in full. It
also runs on the custom Timur-RV32IMC core. AArch64 then changes the instruction set and the word size
(LP64) at the same time and brings a richer exception and memory model. That is the larger step, which is why it comes
second, and it prepares the ground for the 64-bit x86 extension of Phase 15.
\end{decisionbox}

\section{The phase map}

The nineteen phases fall into five stages. Arrows show the build order. Each phase depends on everything before it.
The first four stages lead to Kagan 1.0; the fifth ports it and leads to Kagan 2.0.

\begin{figure}[H]
\centering
\begin{tikzpicture}[
  ph/.style={draw, rounded corners=3pt, line width=0.7pt, minimum width=3.45cm, minimum height=1.3cm,
    text width=3.2cm, align=center, font=\sffamily\scriptsize, inner sep=3pt},
  s1/.style={ph, draw=tkRed!80!black, fill=tkRed!7},
  s2/.style={ph, draw=tkAmber!90!black, fill=tkAmber!9},
  s3/.style={ph, draw=tkBlue!80!black, fill=tkBlue!7},
  s4/.style={ph, draw=tkGreen!80!black, fill=tkGreen!8},
  s5/.style={ph, draw=tkViolet!80!black, fill=tkViolet!8},
  rel/.style={draw=tkViolet!80!black, dashed, rounded corners=3pt, line width=0.7pt, minimum width=2.2cm,
    minimum height=1.3cm, align=center, font=\sffamily\scriptsize, text=tkViolet!80!black, inner sep=3pt},
  flow/.style={-{Stealth[length=2mm]}, line width=0.8pt, draw=tkNavy!70},
]
\node[s1] (p0) at (0,0)      {\textbf{P0\enspace Foundations}\\toolchain, C for kernels};
\node[s1] (p1) at (4.1,0)    {\textbf{P1\enspace Assembly \& boot}\\NASM, BIOS, Multiboot};
\node[s1] (p2) at (8.2,0)    {\textbf{P2\enspace Getting to C}\\VGA, serial, kprintf};
\node[s1] (p3) at (12.3,0)   {\textbf{P3\enspace GDT \& IDT}\\CPU exceptions};
\node[s1] (p4) at (12.3,-2)  {\textbf{P4\enspace Hardware IRQs}\\PIC, timer, keyboard};
\node[s2] (p5) at (8.2,-2)   {\textbf{P5\enspace Physical memory}\\page-frame allocator};
\node[s2] (p6) at (4.1,-2)   {\textbf{P6\enspace Paging}\\higher-half kernel};
\node[s2] (p7) at (0,-2)     {\textbf{P7\enspace Kernel heap}\\kmalloc, lists, tests};
\node[s3] (p8) at (0,-4)     {\textbf{P8\enspace Processes}\\context switch, scheduler};
\node[s3] (p9) at (4.1,-4)   {\textbf{P9\enspace Synchronization}\\locks, semaphores};
\node[s3] (p10) at (8.2,-4)  {\textbf{P10\enspace User mode}\\ring 3, system calls};
\node[s3] (p11) at (12.3,-4) {\textbf{P11\enspace Process API}\\fork, exec, wait, COW};
\node[s4] (p12) at (12.3,-6) {\textbf{P12\enspace Storage}\\ATA driver, buffer cache};
\node[s4] (p13) at (8.2,-6)  {\textbf{P13\enspace File systems}\\VFS, initrd, KaganFS};
\node[s4] (p14) at (4.1,-6)  {\textbf{P14\enspace Userland}\\libc, shell, pipes};
\node[s4] (p15) at (0,-6)    {\textbf{P15\enspace Kagan 1.0}\\hardening, release};
\node[s5] (p16) at (0,-8)    {\textbf{P16\enspace Portability}\\arch layer, devicetree};
\node[s5] (p17) at (4.1,-8)  {\textbf{P17\enspace RISC-V port}\\Sv32, SBI, virtio};
\node[s5] (p18) at (8.2,-8)  {\textbf{P18\enspace ARM port}\\AArch64, GIC, MMU};
\node[rel] (v2) at (11.85,-8) {\textbf{v2.0}\\three architectures};
\foreach \a/\b in {p0/p1,p1/p2,p2/p3,p3/p4,p4/p5,p5/p6,p6/p7,p7/p8,p8/p9,p9/p10,p10/p11,p11/p12,p12/p13,p13/p14,p14/p15,p15/p16,p16/p17,p17/p18}
  \draw[flow] (\a) -- (\b);
\draw[flow, draw=tkViolet!80!black, dashed] (p18) -- (v2);
\end{tikzpicture}

\vspace{0.6em}
{\sffamily\scriptsize
\begin{tabular}{@{}l@{\hspace{3em}}l@{}}
\tikz\fill[tkRed!60] (0,0) rectangle (0.25,0.25); Stage I: Bare metal &
\tikz\fill[tkAmber!70] (0,0) rectangle (0.25,0.25); Stage II: Memory \\[2pt]
\tikz\fill[tkBlue!60] (0,0) rectangle (0.25,0.25); Stage III: Processes and protection &
\tikz\fill[tkGreen!60] (0,0) rectangle (0.25,0.25); Stage IV: Persistence and userland (release 1.0) \\[2pt]
\tikz\fill[tkViolet!60] (0,0) rectangle (0.25,0.25); Stage V: Ports (release 2.0) & \\
\end{tabular}}
\caption{The nineteen phases, grouped into five stages.}
\end{figure}

\textbf{Stage I, Bare metal (Phases 0--4).} The toolchain, the boot path and the first hardware support.
By the end the kernel runs C code, prints to the screen and the serial port, survives CPU exceptions with a readable
error report, and reacts to the timer and the keyboard through a small built-in command monitor.

\textbf{Stage II, Memory (Phases 5--7).} The kernel learns which RAM exists and hands it out page by page, turns on
paging, moves itself to the top gigabyte of the address space, and gets a \texttt{kmalloc} heap. The stage also
establishes the automated test harness and continuous integration used for the rest of the project.

\textbf{Stage III, Processes and protection (Phases 8--11).} Multitasking arrives: kernel threads, context switches,
a preemptive scheduler, then the locks that make concurrency safe. Programs move to ring 3 with their own address
spaces and talk to the kernel through system calls. The stage ends with the UNIX process model: \texttt{fork} with
copy-on-write, \texttt{exec}, \texttt{wait} and \texttt{exit}.

\textbf{Stage IV, Persistence and userland (Phases 12--15).} A disk driver and buffer cache provide storage, a VFS
and KaganFS provide files, and a C library, shell and utilities complete a usable system. The last phase hardens,
measures and documents it, and tags release 1.0; it also specifies the optional extensions.

\textbf{Stage V, Ports (Phases 16--18).} First the architecture layer is separated: everything x86 moves behind
one interface, the generic code learns to compile for other machines, and the kernel learns to read a devicetree
instead of Multiboot (16). Then Kagan is ported to 32-bit RISC-V on QEMU's \texttt{virt} board, with OpenSBI as
the firmware, and optionally to the custom Timur-RV32IMC core (17). Finally it is ported to 64-bit ARM (AArch64) on
QEMU's \texttt{virt} board, and optionally to a Raspberry Pi 4. The stage ends with Kagan 2.0: the same kernel,
userland and KaganFS image on three architectures (18).

\section{Phase summary}

\begin{center}
\small\renewcommand{\arraystretch}{1.35}\setlength{\tabcolsep}{4pt}
\begin{longtable}{@{}c >{\raggedright\arraybackslash}p{0.26\linewidth} >{\raggedright\arraybackslash}p{0.6\linewidth}@{}}
\toprule
\tkth{Phase} & \tkth{Title} & \tkth{Deliverables} \\
\midrule
\endfirsthead
\toprule
\tkth{Phase} & \tkth{Title} & \tkth{Deliverables} \\
\midrule
\endhead
\bottomrule
\endlastfoot
0  & Foundations and toolchain & Cross compiler, QEMU, GDB and a Git repo; C warm-up library with unit tests; verified toolchain sources; threat-model stub \\
1  & x86 assembly and booting & A 512-byte boot sector; a Multiboot kernel that GRUB and QEMU boot \\
2  & Getting to C: screen, serial, kprintf & VGA text driver, serial logging, \texttt{kprintf}, \texttt{panic} \\
3  & Segmentation, IDT and exceptions & GDT and IDT; register dump on every CPU exception \\
4  & Hardware interrupts & PIC, 100\,Hz timer, keyboard driver, kernel command monitor \\
5  & Physical memory management & Memory map parsing, bitmap then buddy page-frame allocator, first kernel tests \\
6  & Paging and the higher-half kernel & Paging on, kernel at \texttt{0xC0000000}, VMM API, page-fault reports, read-only kernel code \\
7  & Kernel heap and data structures & \texttt{kmalloc}/\texttt{kfree} on slab caches, intrusive lists, \texttt{make test} in headless QEMU, CI and a trap-on-UB build \\
8  & Processes and scheduling & Kernel threads, context switch, preemptive round-robin (optional fair scheduler), sleep, FPU state per task \\
9  & Synchronization & Spinlocks, wait queues, mutexes, semaphores, \texttt{refcount\_t}, lock-order checking \\
10 & User mode and system calls & TSS, ring 3, \texttt{int 0x80}, checked user-memory access, validating ELF loader, threat model, first user programs \\
11 & Process API & \texttt{fork} (copy-on-write), \texttt{exec}, \texttt{wait}, \texttt{exit}, \texttt{brk}, basic signals, system call fuzzer \\
12 & Storage & ATA PIO driver, block layer, MBR partitions, buffer cache \\
13 & File systems & VFS, file descriptors, tar initrd, writable KaganFS that rejects corrupted images \\
14 & Userland & Kagan libc, pipes, tty line discipline, \texttt{kash} shell, coreutils \\
15 & Hardening and release & Final security audit, users, permissions and sessions, benchmarks, docs, a reproducible and signed v1.0; optional tracks \\
\midrule
16 & Portability: an architecture layer & x86 moved to \texttt{kernel/arch/i386}, \texttt{include/arch.h}, \texttt{struct bootinfo}, a host-tested devicetree parser, one Makefile for three \texttt{ARCH} values \\
17 & The RISC-V port & Kagan on \texttt{qemu-system-riscv32 -machine virt}: OpenSBI, Sv32, PLIC, virtio-blk, boot to \texttt{kash}; optionally the Timur-RV32IMC core \\
18 & The ARM port & Kagan on \texttt{qemu-system-aarch64 -machine virt}: EL1, GICv2, generic timer, MMU, LP64 userland; v2.0 on three architectures; optionally a Raspberry Pi 4 \\
\end{longtable}
\end{center}



\section{What Kagan 1.0 and 2.0 look like}

Figure~\ref{fig:arch} shows the finished 1.0 system and the phase that builds each part. It is worth coming back to
this picture as the map of where each piece fits. Figure~\ref{fig:arch2} shows what Stage V changes: the
generic kernel stays one body of code, and everything that depends on the processor or the board moves below
\texttt{include/arch.h}, with one backend per architecture.

\begin{figure}[htbp]
\centering
\begin{tikzpicture}[
  layer/.style={draw=tkNavy!60, rounded corners=2pt, minimum height=1.05cm, font=\sffamily\scriptsize,
    align=center, inner sep=2pt, line width=0.6pt, anchor=north west},
  ph/.style={text=tkRed, font=\sffamily\scriptsize\bfseries},
]
\node[layer, fill=tkGreen!10, minimum width=4.2cm, text width=4.0cm] at (0,0)    {init, \texttt{kash} shell, coreutils\\{\color{tkRed}\bfseries P11, P14}};
\node[layer, fill=tkGreen!10, minimum width=4.2cm, text width=4.0cm] at (4.35,0) {Kagan libc: stdio, malloc, string\\{\color{tkRed}\bfseries P14}};
\node[layer, fill=tkGreen!10, minimum width=4.3cm, text width=4.1cm] at (8.7,0)  {user ELF programs\\{\color{tkRed}\bfseries P10, P11}};
\node[layer, fill=tkRed!8, minimum width=13cm, text width=12.8cm, minimum height=0.8cm] at (0,-1.25)
  {System call interface: \texttt{int 0x80}, argument checks, \texttt{copy\_from\_user} / \texttt{copy\_to\_user}\quad{\color{tkRed}\bfseries P10}};
\node[layer, fill=tkBlue!7, minimum width=3.1cm, text width=2.95cm, minimum height=1.25cm] at (0,-2.25)    {Processes and scheduler\\fork, exec, signals\\{\color{tkRed}\bfseries P8, P11}};
\node[layer, fill=tkBlue!7, minimum width=3.1cm, text width=2.95cm, minimum height=1.25cm] at (3.3,-2.25)  {Memory\\PMM, VMM, heap, COW\\{\color{tkRed}\bfseries P5--P7, P11}};
\node[layer, fill=tkBlue!7, minimum width=3.1cm, text width=2.95cm, minimum height=1.25cm] at (6.6,-2.25)  {VFS\\tarfs, KaganFS, devfs, pipes\\{\color{tkRed}\bfseries P13, P14}};
\node[layer, fill=tkBlue!7, minimum width=3.1cm, text width=2.95cm, minimum height=1.25cm] at (9.9,-2.25)  {Synchronization\\spinlocks, mutexes, sems\\{\color{tkRed}\bfseries P9}};
\node[layer, fill=tkAmber!10, minimum width=13cm, text width=12.8cm, minimum height=0.8cm] at (0,-3.7)
  {Drivers: VGA, serial, keyboard, PIT, tty, ATA, RAM disk, block layer, buffer cache\quad{\color{tkRed}\bfseries P2, P4, P12, P14}};
\node[layer, fill=tkAmber!10, minimum width=13cm, text width=12.8cm, minimum height=0.8cm] at (0,-4.7)
  {i386 layer: Multiboot entry, GDT, IDT and ISR stubs, PIC, TSS, paging, context switch\quad{\color{tkRed}\bfseries P1--P4, P6, P8, P10}};
\node[layer, fill=tkGray!12, minimum width=13cm, text width=12.8cm, minimum height=0.8cm] at (0,-5.9)
  {PC hardware (QEMU): CPU, RAM, VGA, 8259 PIC, 8254 PIT, 8042 keyboard controller, IDE disk, 16550 UART};
\draw[decorate, decoration={brace, amplitude=5pt, mirror}, tkGray, line width=0.7pt] (-0.15,0) -- (-0.15,-1.05)
  node[midway, left=6pt, font=\sffamily\scriptsize\bfseries, text=tkGray, align=right] {user\\space\\(ring 3)};
\draw[decorate, decoration={brace, amplitude=5pt, mirror}, tkGray, line width=0.7pt] (-0.15,-1.25) -- (-0.15,-5.5)
  node[midway, left=6pt, font=\sffamily\scriptsize\bfseries, text=tkGray, align=right] {kernel\\(ring 0)};
\draw[decorate, decoration={brace, amplitude=5pt, mirror}, tkGray, line width=0.7pt] (-0.15,-5.9) -- (-0.15,-6.7)
  node[midway, left=6pt, font=\sffamily\scriptsize\bfseries, text=tkGray, align=right] {hardware};
\end{tikzpicture}
\caption{Layers of Kagan 1.0, with the phase that builds each part.}
\label{fig:arch}
\end{figure}

\begin{figure}[htbp]
\centering
\begin{tikzpicture}[
  layer/.style={draw=tkNavy!60, rounded corners=2pt, font=\sffamily\scriptsize, align=center, inner sep=2pt,
    line width=0.6pt, anchor=north west, execute at begin node={\hyphenpenalty=10000\exhyphenpenalty=10000\relax}},
  wide/.style={layer, minimum width=13cm, text width=12.8cm, minimum height=0.9cm},
  col/.style={layer, minimum width=4.2cm, text width=4.0cm, minimum height=1.6cm},
]
\node[wide, fill=tkGreen!10] at (0,0)
  {\textbf{User space}, rebuilt for each architecture: init, \texttt{kash}, coreutils, Kagan libc\\\texttt{crt0} and system-call stubs in \texttt{user/arch/<arch>/}\quad{\color{tkRed}\bfseries P17, P18}};
\node[wide, fill=tkBlue!7] at (0,-1.05)
  {\textbf{Generic kernel}, one source for all three architectures\\system calls, processes and scheduler, memory, VFS, KaganFS, buffer cache, tty, \texttt{ktest}\quad{\color{tkRed}\bfseries P16}};
\node[wide, fill=tkAmber!10] at (0,-2.1)
  {\textbf{Drivers}: the 16550 core with port-I/O or MMIO accessors, ATA (x86), PL011 (ARM)\\virtio-mmio and virtio-blk, shared by both ports; the devicetree parser \texttt{lib/fdt.c}\quad{\color{tkRed}\bfseries P16--P18}};
\node[wide, fill=tkRed!8] at (0,-3.15)
  {\textbf{Architecture interface}: \texttt{include/arch.h} and the \texttt{<arch/*.h>} headers\\\texttt{struct trapframe}, \texttt{struct bootinfo}, \texttt{arch\_switch}, \texttt{arch\_mmu\_*}, \texttt{arch\_irq\_*}\quad{\color{tkRed}\bfseries P16}};
\node[col, fill=tkLight] at (0,-4.25)
  {\textbf{\texttt{kernel/arch/i386}}\\Multiboot entry, GDT, IDT\\8259 PIC, PIT, two-level paging\\\texttt{int 0x80}\quad{\color{tkRed}\bfseries P1--P10, P16}};
\node[col, fill=tkViolet!8, draw=tkViolet!70] at (4.4,-4.25)
  {\textbf{\texttt{kernel/arch/riscv32}}\\entry from OpenSBI, \texttt{stvec} traps\\PLIC, SBI timer, Sv32 paging\\\texttt{ecall}\quad{\color{tkRed}\bfseries P17}};
\node[col, fill=tkViolet!8, draw=tkViolet!70] at (8.8,-4.25)
  {\textbf{\texttt{kernel/arch/aarch64}}\\EL1 entry, \texttt{VBAR\_EL1} vectors\\GICv2, generic timer, MMU\\\texttt{svc \#0}\quad{\color{tkRed}\bfseries P18}};
\node[col, fill=tkGray!12] at (0,-6.0)
  {PC in \texttt{qemu-system-i386}\\8259 PIC, 8254 PIT, IDE disk\\16550 UART by port I/O\\{\color{tkGray}(or a real PC)}};
\node[col, fill=tkGray!12] at (4.4,-6.0)
  {QEMU riscv32 \texttt{virt}\\PLIC, virtio-mmio\\16550 UART at \texttt{0x10000000}\\{\color{tkGray}(optionally Timur-RV32IMC)}};
\node[col, fill=tkGray!12] at (8.8,-6.0)
  {QEMU aarch64 \texttt{virt}, Cortex-A72\\GICv2, virtio-mmio\\PL011 UART\\{\color{tkGray}(optionally a Raspberry Pi 4)}};
\draw[decorate, decoration={brace, amplitude=5pt, mirror}, tkGray, line width=0.7pt] (-0.15,0) -- (-0.15,-0.9)
  node[midway, left=6pt, font=\sffamily\scriptsize\bfseries, text=tkGray, align=right] {user\\space};
\draw[decorate, decoration={brace, amplitude=5pt, mirror}, tkGray, line width=0.7pt] (-0.15,-1.05) -- (-0.15,-4.05)
  node[midway, left=6pt, font=\sffamily\scriptsize\bfseries, text=tkGray, align=right] {generic\\kernel};
\draw[decorate, decoration={brace, amplitude=5pt, mirror}, tkGray, line width=0.7pt] (-0.15,-4.25) -- (-0.15,-5.85)
  node[midway, left=6pt, font=\sffamily\scriptsize\bfseries, text=tkGray, align=right] {one\\backend\\per arch};
\draw[decorate, decoration={brace, amplitude=5pt, mirror}, tkGray, line width=0.7pt] (-0.15,-6.0) -- (-0.15,-7.6)
  node[midway, left=6pt, font=\sffamily\scriptsize\bfseries, text=tkGray, align=right] {boards};
\end{tikzpicture}
\caption{Kagan 2.0 after Stage V: one generic kernel, one architecture interface, three backends and three boards.}
\label{fig:arch2}
\end{figure}

The source tree grows into the layout below; each phase document states which directories it adds. The entries marked \texttt{(Stage V)} arrive in Phases 16--18.

\begin{lst}{plain}{kagan-os/ (layout at 1.0, plus the Stage V additions)}
kagan-os/
|-- docs/                  design documents (LaTeX), journal.md, security/
|   `-- porting.md         the leak audit and architecture-interface notes        (Stage V)
|-- boot/                  loader.s (Multiboot entry), grub.cfg
|-- kernel/
|   |-- arch/i386/         gdt.c idt.c isr.s irq.c pic.c pit.c tss.c paging.c switch.s
|   |   `-- include/arch/  the x86 headers behind the architecture interface      (Stage V)
|   |-- arch/riscv32/      boot.S entry.S trap.c sbi.c plic.c timer.c sv32.c      (Stage V)
|   |                      switch.S user.S arch.mk kernel.ld include/arch/
|   |-- arch/aarch64/      boot.S vectors.S trap.c gic.c timer.c mmu.c            (Stage V)
|   |                      switch.S user.S arch.mk kernel.ld include/arch/
|   |-- mm/                pmm.c vmm.c heap.c
|   |-- proc/              task.c sched.c fork.c exec.c signal.c
|   |-- sync/              spinlock.c mutex.c semaphore.c waitq.c
|   |-- fs/                vfs.c file.c tarfs.c kaganfs.c devfs.c pipe.c
|   |-- drivers/           vga.c serial.c keyboard.c tty.c ata.c ramdisk.c blockdev.c bcache.c
|   |                      + uart16550.c pl011.c virtio_mmio.c virtio_blk.c       (Stage V)
|   |-- syscall/           syscall.c uaccess.c
|   `-- kmain.c
|-- lib/                   libk: string.c printf.c bitmap.c ringbuf.c list.h
|                          + fdt.c, the devicetree parser                         (Stage V)
|-- include/               kernel headers
|                          + arch.h and bootinfo.h                                (Stage V)
|-- user/
|   |-- libc/              Kagan libc
|   |-- arch/<arch>/       crt0 and system call stubs for each architecture       (Stage V)
|   |-- init/ sh/ bin/     init, kash, coreutils
|   `-- user.ld
|-- tests/                 ktest suites
|   `-- data/              devicetree blobs dumped from QEMU, for the FDT tests   (Stage V)
|-- tools/                 mkkaganfs.c and helper scripts
|-- link.ld
`-- Makefile
\end{lst}

\section{Risks and mitigations}

Kernel development has predictable failure modes, the worst of them a machine that silently reboots with no output.
The table lists the project risks and their mitigations; the protocol after it is the standard debugging procedure.

\begin{center}
\small\renewcommand{\arraystretch}{1.35}\setlength{\tabcolsep}{4pt}
\begin{longtable}{@{}>{\raggedright\arraybackslash}p{0.33\linewidth} c >{\raggedright\arraybackslash}p{0.5\linewidth}@{}}
\toprule
\tkth{Risk} & \tkth{Likelihood} & \tkth{Mitigation} \\
\midrule
\endfirsthead
\toprule
\tkth{Risk} & \tkth{Likelihood} & \tkth{Mitigation} \\
\midrule
\endhead
\bottomrule
\endlastfoot
Triple faults (instant reboot) with no output & High & Serial logging from Phase 2 on; run QEMU with \texttt{-d int,cpu\_reset -no-reboot}; GDB; small commits so \texttt{git bisect} can find the change that broke it. \\
Errors in hand-written assembly & High & Keep assembly to small stubs behind C interfaces, check it against compiler output, and single-step every stub once in GDB. \\
Memory corruption that shows up far from its cause & High & Assertions everywhere, magic numbers in headers, guard pages, \texttt{heap\_check()}, deterministic tests. \\
A single defect blocks progress for days & High & Follow the debugging protocol below; time-box the investigation, then change approach. \\
Scope creep (GUI or networking too early) & Medium & Exit criteria are the definition of done. New ideas go on the extensions list in Phase 15. \\
A user program or a corrupted disk image can crash the kernel, found only at the end & Medium & Each phase that adds an entry point adds its negative tests at once: evil programs from Phase 10, the system call fuzzer from Phase 11, host fuzzing of every parser from Phase 13, all run in CI. \\
Toolchain trouble (cross-GCC build, Flatpak sandbox) & Medium & Phase 0 checklist, a host terminal, documented fallback compiler flags. \\
Outdated tooling in \LOB & High & The phase docs give current equivalents (\texttt{grub2-mkrescue}, QEMU); the OSDev Wiki has errata. \\
Long single-maintainer project stalls & Medium & Every phase ends in a testable, demonstrable milestone; progress and decisions are recorded in the engineering journal. \\
\midrule
Stage V: x86 habits carried into the ports (a forgotten \texttt{sfence.vma}, missing TLB maintenance or barriers) & High & Check every step against the specification, single-step in GDB, and compare with xv6-riscv or Linux. \\
Stage V: QEMU is more forgiving than real hardware (caches, alignment, uninitialised registers) & Medium & Keep real hardware optional and last: the QEMU milestone comes first and counts on its own. \\
Stage V: toolchain trouble with new targets (multilib, \texttt{-march} strings with \texttt{\_zicsr}) & Medium & Section~9 of \texttt{docs/setup.md}, \emph{Toolchains for the ports}, and the Phase 16 checklist (Task 16.1). \\
Stage V: the custom core lacks a feature Kagan needs (S-mode, an MMU, a timer) & Medium & Task 17.13's requirements checklist. The QEMU milestone stands on its own without the core. \\
Stage V: scope creep (SMP, graphics on the ports) & Medium & The Stage V exit criteria are the definition of done. \\
\end{longtable}
\end{center}

\begin{tipbox}[Debugging protocol]
\begin{enumerate}
  \item \textbf{Reproduce it reliably}, then make it smaller: remove code until the shortest path to the bug remains.
  \item \textbf{Look, don't guess.} Use the serial log, \texttt{-d int}, the QEMU monitor (\texttt{info registers}, \texttt{info mem}), and GDB breakpoints.
  \item \textbf{Write the hypothesis down} in the engineering journal before changing code. Change one thing at a time.
  \item \textbf{Re-read the section} of the Intel manual (in Stage V, the RISC-V or Arm specification) or the reference for the exact hardware rule involved.
  \item \textbf{Compare with MINIX} (\OSDI source) or xv6: how does a working kernel do the same step?
  \item \textbf{After two sessions without progress}, escalate: write the problem up with the journal notes and ask on the OSDev forum or a relevant mailing list. Writing it up often exposes the cause.
\end{enumerate}
\end{tipbox}

\section{Exit criteria}

A phase is complete when every one of its exit criteria holds.

\begin{checklist}
  \item \textbf{P0}\enspace \texttt{i686-elf-gcc}, NASM, QEMU and GDB work; repo created; toolchain signatures verified; C warm-up library and overflow helpers pass their tests; threat-model stub written.
  \item \textbf{P1}\enspace Boot sector prints a message; Multiboot kernel boots via \texttt{-kernel} and via a GRUB ISO.
  \item \textbf{P2}\enspace Colored banner on screen, same log on the serial port, \texttt{kprintf} and \texttt{panic} tested; the kernel builds with \texttt{-Werror}.
  \item \textbf{P3}\enspace Own GDT and IDT loaded; divide-by-zero shows a register dump; \texttt{int3} returns cleanly.
  \item \textbf{P4}\enspace Timer ticks at 100\,Hz; keyboard input with shift and caps; monitor commands run.
  \item \textbf{P5}\enspace Memory map printed; bitmap and then buddy allocator pass the tests from 16\,MiB to 3\,GiB of RAM.
  \item \textbf{P6}\enspace Kernel runs at \texttt{0xC0000000}; null dereference produces a clean page-fault report; a write to kernel code faults.
  \item \textbf{P7}\enspace 10\,000-operation stress test passes through the slab caches; \texttt{make test} runs headless and in CI; the trap-on-UB build passes; allocation failures are handled.
  \item \textbf{P8}\enspace Several kernel threads preempt each other; \texttt{sleep} works; no leaks over 1000 task cycles; FPU state stays with its task.
  \item \textbf{P9}\enspace Race demo fixed by locks; bounded buffer and dining philosophers run without deadlock; lock order checked.
  \item \textbf{P10}\enspace Ring-3 program prints through \texttt{write}; faulty and malicious programs are killed or get errors, the kernel survives; malformed ELF files are rejected; every system call has a row in the boundary table.
  \item \textbf{P11}\enspace \texttt{fork}/\texttt{exec}/\texttt{wait} chain works; copy-on-write verified, also through \texttt{copy\_to\_user}; no leaks over 1000 forks; ten minutes of fuzzing without a panic.
  \item \textbf{P12}\enspace QEMU disk identified; data survives reboot; partitions detected and corrupted tables rejected; buffer cache in use.
  \item \textbf{P13}\enspace Files written by Kagan pass \texttt{fsck.minix} on the host; file descriptor semantics tested; corrupted images are rejected without a panic; the parsers survive an hour of host fuzzing.
  \item \textbf{P14}\enspace Boots to the \texttt{kash} prompt; pipes, redirection and Ctrl+C work; coreutils present; resource limits enforced.
  \item \textbf{P15}\enspace One hour of fuzzing as an ordinary user without a panic; permission and session tests pass; benchmarks, docs and the threat model finished; two clean builds match and v1.0 is tagged and signed.
  \item \textbf{P16}\enspace Generic code free of x86 (audit in \texttt{docs/porting.md}); \texttt{make ARCH=i386 test} passes; generic code compiles for riscv32 and aarch64 without warnings; FDT parser passes its host tests.
  \item \textbf{P17}\enspace \texttt{qemu-system-riscv32 -machine virt} boots to the \texttt{kash} prompt; \texttt{make ARCH=riscv32 test} passes, evil programs included; W\textasciicircum X holds; a KaganFS image written on RISC-V passes \texttt{fsck.minix}.
  \item \textbf{P18}\enspace \texttt{qemu-system-aarch64 -machine virt} boots to the \texttt{kash} prompt; \texttt{make ARCH=aarch64 test} passes, evil programs included; W\textasciicircum X holds; the same KaganFS image works on all three; v2.0 tagged.
\end{checklist}

\booklegend
\portlegend
\end{document}
